% Chapter 1

\chapter{Introducción General}
\label{Chapter1}

Este capítulo introduce el problema de la predicción de series de tiempo financieras y su aplicación en un sistema de 	\textit{trading} de criptomonedas. Se describe el contexto y la motivación para el uso de la arquitectura 	\textit{Transformer}, se presenta la problemática del 	\textit{trading} algorítmico, se revisa el estado del arte y, finalmente, se definen los objetivos y el alcance del trabajo.

\section{Motivación}

La elección de este trabajo se fundamenta en dos pilares: el problema a resolver y las herramientas seleccionadas para abordarlo. El 	\textit{trading} de activos financieros, una actividad que busca generar rentabilidad a partir de las fluctuaciones de precios, representa un desafío complejo y dinámico. Gran parte de las decisiones en este campo se basan en el "análisis técnico", que consiste en la interpretación de indicadores derivados de precios y volúmenes pasados. La combinación manual de estos indicadores para tomar decisiones fundamentadas requiere una experiencia considerable y es difícil de escalar para analizar múltiples indicadores simultáneamente. Esto hace atractiva la idea de desarrollar un modelo de 	\textit{machine learning} que pueda procesar una gran cantidad de indicadores técnicos para optimizar las decisiones de compra y venta en cada instante.

Por otro lado, la arquitectura 	\textit{Transformer}, gracias a su mecanismo de auto-atención, ha demostrado una capacidad excepcional para modelar dependencias a largo plazo en secuencias, consolidándose como el estado del arte en el procesamiento de lenguaje natural. Esta cualidad la convierte en una candidata ideal para el análisis de series de tiempo financieras, donde eventos pasados pueden tener una influencia significativa en el futuro. A diferencia de los modelos recurrentes, el 	\textit{Transformer} puede teóricamente capturar relaciones entre puntos distantes en el tiempo, permitiendo evitar operaciones rentables pero de alto riesgo o, por el contrario, aprovechar oportunidades de baja rentabilidad con un riesgo despreciable.

\section{Descripción de la Problemática}

En la comercialización de activos financieros, el objetivo es generar una ganancia comprando a un precio bajo y vendiendo a un precio más alto. Esto sería trivial si se pudiera predecir el futuro de los precios con certeza, pero la realidad de los mercados es inherentemente estocástica, como se puede apreciar en la Figura \ref{fig:btc_price_chart}. Quienes se dedican a esta actividad pueden clasificarse a grandes rasgos en dos categorías: analistas fundamentales, que basan sus decisiones en el valor intrínseco de los activos, y analistas técnicos, que se centran en el estudio de patrones históricos de precios y volúmenes.

\begin{figure}[h!]
    \centering
    \includegraphics[width=\textwidth]{"Figures/Screenshot From 2026-03-17 21-57-54.png"}
    \caption{Evolución del precio del par BTC/USDT en el período analizado.\protect\footnotemark}
    \label{fig:btc_price_chart}
\end{figure}
\footnotetext{Gráfico obtenido de la plataforma de trading Binance.}

Este trabajo se enmarca en el análisis técnico, buscando desarrollar un sistema de 	\textit{trading} algorítmico para operar de manera autónoma en el mercado 	\textit{spot} de criptomonedas. Específicamente, se trabajará con el par BTC/USDT, donde el objetivo es decidir en momentos discretos si es conveniente mantener la posición en Bitcoin (esperando que su valor aumente respecto al Tether) o en USDT (protegiéndose de una posible caída del Bitcoin).

\section{Estado del Arte}

El uso de inteligencia artificial para el 	\textit{trading} ha ganado un protagonismo creciente en las últimas dos décadas, tanto en el ámbito comercial como en el académico. Empresas como 	\textit{AutoTrade.AI} o 	\textit{Finbrain Technologies} ya ofrecen servicios de 	\textit{trading} automático basados en señales predictivas generadas por modelos de IA.

El ámbito académico ha explorado extensamente el uso de 	\textit{machine learning} para la optimización de portafolios y la predicción de mercados. Estudios recientes han aplicado diversos modelos para predecir los movimientos a corto plazo de criptomonedas como Bitcoin, utilizando variables de análisis técnico, datos de la 	\textit{blockchain} y análisis de sentimiento del mercado. Sin embargo, muchos de estos enfoques se enfrentan a un obstáculo fundamental: la Hipótesis de los Mercados Eficientes (	\textit{EMH}). Esta hipótesis sugiere que toda la información disponible ya está reflejada en el precio de los activos, lo que haría inútil cualquier intento de obtener ganancias extraordinarias. Aunque la validez de la 	\textit{EMH} para el mercado de criptomonedas es un tema de debate, la eficiencia del mercado parece aumentar a medida que este madura, lo que representa un desafío constante para el desarrollo de modelos predictivos.

\section{El Desafío de los Datos Financieros}

La aplicación de 	\textit{machine learning} a los mercados financieros presenta desafíos únicos:

\begin{itemize}
    \item \textbf{Bajo ratio señal-ruido:} Los movimientos de precios son inherentemente ruidosos. La "señal" predictiva a menudo se encuentra oculta bajo una gran cantidad de fluctuaciones aleatorias o "ruido", lo que incrementa el riesgo de sobreajuste.

    \item \textbf{No estacionariedad:} Las propiedades estadísticas de las series de tiempo financieras cambian con el tiempo, haciendo que los patrones aprendidos en un período puedan no ser válidos en el futuro.

    \item \textbf{Costos de transacción:} Operar con alta frecuencia, aunque teóricamente podría capturar más oportunidades, puede llevar a que las ganancias sean consumidas por los costos transaccionales.
\end{itemize}

Estos factores subrayan la necesidad de un diseño metodológico riguroso y específico para el dominio financiero.

\section{Objetivos, Alcance y Requerimientos}

El objetivo principal de este trabajo es estudiar y evaluar rigurosamente el desempeño de la arquitectura 	\textit{Transformer} en la predicción de series de tiempo para su aplicación en un sistema de 	\textit{trading} de criptomonedas en el mercado 	\textit{spot} de la plataforma 	\textit{Binance}.

Los objetivos específicos son:
\begin{enumerate}
    \item Realizar una revisión del estado del arte sobre la aplicación de modelos de 	\textit{deep learning}, y en particular 	\textit{Transformer}, a la predicción de series de tiempo.
    \item Implementar un modelo 	\textit{Transformer} para estimar si es óptimo posicionarse en BTC o USDT en el par BTC/USDT.
    \item Comparar el desempeño del modelo con una estrategia de referencia de "comprar y esperar" (	\textit{buy and hold}). El objetivo final no es la cantidad de USDT acumulado, sino la cantidad de BTC.
    \item Probar distintas técnicas y comparar los resultados obtenidos para decidir el modelo final.
\end{enumerate}

El alcance de este trabajo se centra en la implementación y evaluación de un sistema de \textit{trading} algorítmico, cuya arquitectura conceptual se ilustra en la Figura \ref{fig:trading_system_architecture}. Este sistema incluye su operación simulada en un entorno de \textit{backtesting} y la capacidad de operar en tiempo real. Adicionalmente, se buscará una frecuencia de operación que optimice el equilibrio entre la disponibilidad de datos para entrenamiento y el impacto de los costos transaccionales.

\begin{figure}[h!]
    \centering
    \includegraphics[width=0.9\textwidth]{"Figures/bloques-tp-final.drawio.png"}
    \caption{Diagrama conceptual de la arquitectura del sistema de trading.}
    \label{fig:trading_system_architecture}
\end{figure}
